\chapter{The Hierarchy of Forces: A Unified Derivation}

\author{James Tyndall}

%----------------------------------------------------------------------
\section{Introduction: From a Single Mechanism to a Diverse Universe}
%----------------------------------------------------------------------

Why is the force that binds a nucleus quintillions of times stronger than the force that holds a planet in orbit?  The Standard Model answers with a patchwork of different fields and coupling constants.  SDT answers with a single mechanism: \textbf{all forces are manifestations of the spation pressure field, and their apparent differences arise from the different geometric ways in which particles interact with this field.}

This chapter provides a first-principles derivation of the four known fundamental forces as a mechanical \textbf{hierarchy of repercussions}.


%----------------------------------------------------------------------
\section{Order 0: Baryonic Confinement (The Strong Force)}
%----------------------------------------------------------------------

The most powerful force is the most direct interaction: the universal background pressure ($P_0$) confining a baryonic vortex.

\textbf{Mechanism:} The proton vortex displaces the spation medium; the immense pressure pushes back, holding the vortex together.  Self-confinement.

\textbf{Derivation:}
\begin{align}
  P_0 &\approx 8.0 \times 10^{34}\;\text{Pa} \quad\text{(from $E = mc^2$ and proton geometry)} \\
  A_{\text{proton}} &= \pi R_p^2 \approx 2.22 \times 10^{-30}\;\text{m}^2
\end{align}

\begin{equation}\label{eq:F-strong}
  \boxed{F_{\text{Strong}} = P_0 \times A_{\text{proton}} \approx 1.77 \times 10^5\;\text{N}}
\end{equation}

This is the baseline, Order-0 interaction: a direct measure of the universe's fundamental pressure.


%----------------------------------------------------------------------
\section{Order 1: Leptonic Confinement (The SDT ``Weak Force'')}
%----------------------------------------------------------------------

\textbf{Mechanism:} The electron is also a self-confining vortex, but its stability is governed by the \textbf{local pressure gradient} created by a nearby nucleus---not $P_0$ directly.

\textbf{Derivation} (electron in ground-state hydrogen):

Pressure of the proton's propagated field at the Bohr radius:
\begin{equation}
  P_{\text{local}} = P_0 \left(\frac{R_p}{r_{\text{Bohr}}}\right)^2
    = 8.0 \times 10^{34} \left(\frac{0.84 \times 10^{-15}}{5.29 \times 10^{-11}}\right)^2
    \approx 2.02 \times 10^{25}\;\text{Pa}
\end{equation}

\begin{equation}\label{eq:F-leptonic}
  F_{\text{Leptonic}} = P_{\text{local}} \times \pi R_e^2
    \approx 2.02 \times 10^{25} \times \pi(2.82 \times 10^{-15})^2
    \approx 5.03 \times 10^{-4}\;\text{N}
\end{equation}

Order-1: pressure from a \emph{propagated} field, $\sim 10^8$ weaker than Order-0.


%----------------------------------------------------------------------
\section{Order 2: The Unified Electro-Gravitational Force}
%----------------------------------------------------------------------

This is the force \emph{between} two particles, arising from the interaction of two propagated pressure fields.

\textbf{Derivation} (proton--electron in hydrogen):
\begin{equation}\label{eq:F-EM}
  F_{\text{interaction}} = \frac{c^2 R_p}{k_p^2} \cdot \frac{m_e}{r_e^2}
    = \frac{(2.998 \times 10^8)^2 \times 0.84 \times 10^{-15}}{(0.546)^2}
      \cdot \frac{9.109 \times 10^{-31}}{(5.29 \times 10^{-11})^2}
    \approx 8.24 \times 10^{-8}\;\text{N}
\end{equation}

This is identical to the \textbf{electrostatic force}---$\sim 10^5$ weaker than Order-1.


%----------------------------------------------------------------------
\section{Order 3: Macroscopic Gravity}
%----------------------------------------------------------------------

The faintest echo: the \textbf{occlusion} of propagated pressure fields.

\textbf{Derivation} (proton--electron Newtonian gravity):
\begin{equation}\label{eq:F-gravity}
  F_{\text{Gravity}} = G\,\frac{m_p\,m_e}{r_e^2} \approx 3.63 \times 10^{-47}\;\text{N}
\end{equation}

Order-3: $\sim 10^{39}$ weaker than the electrostatic force.  This is a third-order repercussion: the occlusion of the propagated leakage of the primary confinement pressure.


%----------------------------------------------------------------------
\section{The Hierarchy Summarised}
%----------------------------------------------------------------------

\begin{center}
\begin{tabular}{llll}
\hline
\textbf{Force (SDT Name)} & \textbf{Order} & \textbf{Mechanism} & \textbf{Strength (N)} \\
\hline
Baryonic Confinement & 0 & Direct $P_0$ self-confinement & $\sim 1.8 \times 10^5$ \\
Leptonic Confinement & 1 & $P_{\text{local}}$ self-confinement & $\sim 5.0 \times 10^{-4}$ \\
Unified Interaction (EM) & 2 & Propagated field interaction & $\sim 8.2 \times 10^{-8}$ \\
Gravitational Occlusion & 3 & Occlusion of propagated field & $\sim 3.6 \times 10^{-47}$ \\
\hline
\end{tabular}
\end{center}


%----------------------------------------------------------------------
\section{The Role of the Weak Force of Decay}
%----------------------------------------------------------------------

The ``Weak Force'' of the Standard Model is not a force in this hierarchy.  It is a measure of \textbf{vortex instability}.  Its ``weakness'' is a measure of the extreme stability of the composite neutron vortex, which takes a long time to overcome its energy barrier and decay.


%----------------------------------------------------------------------
\section{Conclusion}
%----------------------------------------------------------------------

The hierarchy of fundamental forces is not a mystery of arbitrary coupling constants.  It is a direct, predictable, mechanical consequence of a tiered system of interactions within a single, unified pressure field.  The Strong Force is the direct pressure of the universe.  The Electromagnetic Force is the first echo.  Gravity is the faintest, second-order echo of that echo.

This model not only explains the relative strengths but demonstrates their profound, underlying unity.
